\documentclass[10pt,a4paper]{article}
\usepackage[margin=24mm]{geometry}
\usepackage{amsmath,amssymb,booktabs,hyperref}
\title{The Geometry of Natality: Conditional Certificates and Falsifiable Creation-Enabled Repair Models}
\author{Research-stage technical manuscript; author metadata to be confirmed}
\date{19 October 2026}
\begin{document}\maketitle
\begin{abstract}
We examine creation-enabled finite repair planning with independently checked optimization algorithms, conservative symmetry and decomposition certificates, and an explicit dynamic rewrite grammar. The first frozen benchmark comprises 72 small synthetic cases and 1,440 paired measured trials; the finite quotient was slower on the tested symmetry cases. A second closed grammar admits relay creation prerequisites, capacity, and negative application conditions and agrees with an external HiGHS mixed-integer solver in 4,096 enumerated models. A third grammar models creation history, action-time exclusion, link rewiring and deletion and compares a native uniform-cost solver with a separate NetworkX shortest-path computation on an independently generated transition graph. Standard equivalence and exchange arguments explain surviving claims; counterexamples invalidate broader claims. Neither a new foundational theorem nor unrestricted graph-rewriting correctness is established.
\end{abstract}
\section{Introduction and research questions}
Creation-enabled repair permits installing resources or vertices not present initially. This creates questions of future-relevant history, resource coupling, name identity, and action ordering. Our principal questions are whether exact symmetry quotienting and decomposition preserve minimum cost, whether richer creation dependencies yield a distinctive theorem, and whether independent solvers agree on the same bounded semantics. The work prioritizes falsification and precise refusal conditions rather than asserting universal gains.
\section{Finite repair planning}
A bounded planning instance is a finite state set $S$, initial state $s_0$, enabled actions $A(s)$, deterministic successor $T(s,a)$, nonnegative integral cost $c(s,a)$, forbidden predicate $F(s)$ and goal $G(s)$. A valid path uses enabled actions, never visits forbidden states and reaches a goal. Its optimum is $C^*=\min_{\pi}\sum_{a\in\pi}c(a)$, with $C^*=\infty$ if no goal is reachable. A budget $B$ additionally requires $C^*\le B$. State keys must retain every variable affecting future transitions or goals.
\section{Symmetry and decomposition certificates}
Let a finite group $\Gamma$ act on states and actions. When enabledness, transitions, costs and goals are equivariant, identifying states in the same orbit preserves optimal cost. The concrete implementation enumerates finite name permutations in a restricted service-and-replica grammar; it refuses identity-specific costs, opaque guards and other uncertified features. These conditions are essential, not cosmetic. For independent component transition systems with additive costs and conjunctive goals, and with no shared allocation, negative cross-component guards or capacity, standard product decomposition yields
\[ C^*=\sum_{i=1}^{n}C_i^*. \]
The lower bound follows by projecting any global plan to its local components. The upper bound follows by concatenating independent local optima. Shared restart capacity invalidates the premise: two components with restart cost 1 and alternative cost 3 have unconstrained total 2, but capacity one raises the optimum to 4.
\section{Prompt 6: falsification of greedy hypotheses}
The cheapest absolute restart is not generally the best restart. With restart/alternative pairs $(1,1)$ and $(2,4)$ and at most one restart, choosing the first costs 5 while choosing the second costs 3. In the independent binary-choice model, writing $d_i=q_i-r_i$, selecting the largest positive savings gives
\[ C^*=\sum_i q_i-\sum_{j=1}^{\min(k,m)}d^+_{(j)}. \]
An exchange argument proves this standard cardinality-constrained greedy formula. It is not claimed as novel; 80,973 finite combinations were checked.
\section{Prompt 11: relay creation as Set Cover}
An initially isolated terminal $t$ becomes reachable from root exactly when at least one selected relay template covers $t$. Selecting templates with nonnegative costs is therefore weighted Set Cover. A direct reduction from Set Cover gives NP-hardness and membership in NP for the restricted decision problem. This is a standard encoding, not a new complexity discovery. Two hypotheses failed: cheapest-cover-first can cost 2 rather than optimum 1, and an additional zero-cost option need not strictly improve a zero optimum. Independent explicit graph traversal and subset-mask Dijkstra agreed on 9,261 deterministic small instances.
\section{Prompt 12: dependencies and negative application conditions}
Let binary variable $x_i$ indicate creation of relay $i$. Acyclic prerequisite relation $j\prec i$ requires $x_i\le x_j$. Terminal coverage, creation-slot capacity $K$, and forbidden pair $(a,b)$ give
\begin{align*}
\sum_{i:t\in S_i}x_i&\ge1 &&\text{for each terminal }t,\\
x_i&\le x_j &&\text{for each prerequisite }j\prec i,\\
\sum_i x_i&\le K,\\
x_a+x_b&\le1 &&\text{for forbidden pairs}.
\end{align*}
Minimize $\sum_i c_i x_i$. In the closed acyclic grammar, every selected relay is root-reachable by induction on the prerequisite ordering; hence terminal coverage is equivalent to graph reachability. Conversely every legal creation sequence yields a feasible binary vector. This establishes exact correspondence with the displayed integer program, but is a conventional encoding. An independently maintained SciPy/HiGHS MILP solver agreed with explicit graph-subset enumeration on 4,096 deterministic instances. Cheapest relay choice can fail because of prerequisite costs, adding a template need not strictly improve optimum, and negative guards destroy upward closure of selected feasible sets.
\section{Prompt 13: dynamic allocation and graph rewriting}
The new bounded dynamic grammar has a finite name pool $N$, a terminal set $U$, a shared live-relay capacity $K$, and state $s=(A,L,H)$ where $A\subseteq N$ are live relay identities, $L\subseteq A\times U$ are directed relay-to-terminal links, and $H\subseteq N$ is the permanently used-name history. Create$(i)$ requires $i\notin H$, $|A|<K$, and satisfaction of the action-time negative pair guard; it updates $A$ and $H$. Link$(i,t)$ inserts $(i,t)$; unlink removes it; delete$(i)$ removes the relay and all incident links but retains $H$. The goal is that every terminal has a link from a live relay. All actions have nonnegative costs. This model implements bounded link rewriting and deletion, not arbitrary inter-relay rewiring or dynamic prerequisite changes.
\paragraph{Falsification.} (H1) Deleting a relay does not permit reusing its name: $H$ records prior allocation. (H2) Two individually valid creations may not coexist under a pairwise negative application condition. (H3) In this restricted monotone terminal-link goal, there exists an optimal solution using only creates and links: any delete or unlink can be removed together with later actions dependent solely on the removed identity; more directly, retain only the live relays and links in a successful final state, create those relays in their original name order, and add their final links. This removes nonnegative-cost operations and cannot violate live capacity or pairwise exclusion because the final live subset was legal. This proof assumes no prerequisites, transient goals, positive side effects of deletion, or action-time guards beyond pairwise exclusion. It is a restricted normalization lemma, not an originality claim.
\paragraph{Independent oracle.} A native Dijkstra search enumerates transitions from $s_0$. A separate enumerator constructs a directed state graph without calling the native action generator; NetworkX computes weighted single-source shortest paths on that graph. Every bounded test compares exact goal costs. This is an external algorithmic library but not an independent GROOVE rule semantics.
\section{Frozen benchmark and descriptive statistics}
The Prompt 7 benchmark freezes eight strata, three seeds and three sizes: 72 synthetic cases, 36 matched comparison cells and 36 refusal controls. Prompt 9's isolated rerun performed three warmups and 20 paired measured repetitions per algorithm per supported cell, for 1,440 measured trials, all uncensored. The limits were 60 seconds per subprocess and 2 GiB virtual address space. The median of nine instance-level paired differences in milliseconds (first minus second) was:
\begin{table}[ht]\centering\caption{Descriptive paired median timing differences (ms).}\begin{tabular}{lr}\toprule Comparison & Isolated rerun\\\midrule Plain minus quotient UCS & -3.851\\ Full minus decomposition: independent & +0.743\\ Full minus decomposition: local history & +0.740\\ Full minus decomposition: local negative & +0.427\\\bottomrule\end{tabular}\end{table}
The negative quotient result means the quotient was slower on the tested small synthetic cases. Repeated timing measurements on one CI host do not provide a representative sampling distribution over real networks.
\section{Limitations and decision}
No unrestricted graph-rewriting interpreter, cyclic prerequisite support, real network traces, independent GROOVE run, complete original-manuscript bibliography audit, or genuinely novel foundational theorem is demonstrated. Prompt 12's external HiGHS MILP baseline applies only to the static closed dependency model. Prompt 13's external NetworkX baseline checks the bounded dynamic state graph, not a separately authored graph grammar. Claims about arbitrary negative application conditions, name recycling or inter-relay rewiring would exceed implemented semantics. Consequently this is a research-stage technical draft, not submission-ready Q1 work.
\section*{Availability and evidence}
Repository: \url{https://github.com/Arithmetic-Power-Geometry/The-Geometry-of-Natality}. Branch: \texttt{research/gon-stage1-20261019}. Prompt 9 run: \url{https://github.com/Arithmetic-Power-Geometry/The-Geometry-of-Natality/actions/runs/39208415337}. Prompt 12 external MILP: \url{https://github.com/Arithmetic-Power-Geometry/The-Geometry-of-Natality/actions/runs/39210823910}.
\section*{Selected prior-art leads; not a fully audited bibliography}
Classical Dijkstra shortest-path search; weighted Set Cover; integer programming formulations for Steiner connectivity; graph transformation optimization; finite permutation group orbit search. These categories require full authoritative citation verification before journal submission. See research audit files for specific DOI leads.
\end{document}